\documentclass[11pt]{article}
\usepackage[T1]{fontenc}
\usepackage[utf8]{inputenc}
\usepackage[spanish]{babel}
\usepackage{amsmath,amssymb,booktabs}
\usepackage[margin=2.3cm]{geometry}
\usepackage[colorlinks=true,urlcolor=blue,linkcolor=blue]{hyperref}
\newcommand{\Tr}{\operatorname{Tr}}
\newcommand{\E}{\mathbb{E}}
\newcommand{\CUE}{\operatorname{CUE}}
\title{Factor de forma espectral de un caminante bidimensional\\
con moneda homogénea CUE(4)}
\author{Nota de cálculo sobre el modelo de \texttt{Chaos\_Results/main.pdf}}
\date{30 de septiembre de 2026}
\begin{document}
\maketitle

\begin{abstract}
Para una DTQW con $N$ posiciones y una moneda $C\sim\CUE(4)$,
la dimensión espectral es $D=4N$. Si el operador completo presenta
correlaciones caóticas de clase unitaria, su referencia universal es
$K(\tau)=\min(\tau,1)$, con $\tau=t/D$.
Esta referencia no es una identidad para $U=S(I_N\otimes C)$:
la moneda homogénea produce correlaciones entre caminos y la reflexión
genera contribuciones geométricas a tiempos cortos. Se derivan la
referencia CUE, los momentos Haar exactos a uno y dos pasos y una
comprobación numérica reproducible en dos tamaños de Sinaí.
\end{abstract}

\section{Modelo y promedio}
La sección I de \texttt{main.pdf}, p.~2, define $U=S\cdot C$ y sitúa las
condiciones de frontera en $S$; la sección III.D, pp.~15--17, muestra el
SFF. La implementación usada aquí se toma de
\texttt{QWMisc.wl} y de los archivos \texttt{Sinai.wl} y
\texttt{Common.wl} del repositorio. En el orden posición por moneda,
\begin{equation}
 \mathcal H=\mathbb C^N\otimes\mathbb C^4,\qquad
 U_C=S(I_N\otimes C),\qquad D=4N,\qquad C\sim\CUE(4).
 \label{eq:model}
\end{equation}
El orden moneda por posición del PDF da la misma expresión tras una
permutación de la base. Se sortea una moneda por realización y se usa
\emph{la misma} en todos los sitios y en todos los pasos. No se introduce
desorden espacial ni se vuelve a sortear la moneda durante la evolución.

Los cuatro estados son arriba, abajo, derecha e izquierda. Si el vecino
en la dirección $b$ pertenece al dominio, $S$ mueve al caminante y
conserva $b$. Si falta el vecino, mantiene la posición e invierte la
dirección $b\mapsto\bar b$. Este $S$ es una permutación unitaria.
El dominio discreto de \texttt{GenerateSinaiBasis} es
\begin{equation}
 \Omega_{L,R}=\{(x,y)\in\mathbb Z^2:0\leq y\leq x\leq L,
                   \ x^2+y^2\geq R^2\}.
\end{equation}
Se trata del dominio de un octavo con su propia reflexión local.
Para una moneda CUE genérica no se puede identificar automáticamente
con un sector de simetría del Sinaí completo: esa identificación requiere
que la moneda respete las simetrías correspondientes.

Definimos el SFF físico sin normalizar, y su versión normalizada, como
\begin{align}
 F_\Omega(t)&=\int_{U(4)}\left|\Tr\,[S(I_N\otimes C)]^t\right|^2
                         d\mu_H(C), \label{eq:exactintegral}\\
 K_\Omega(t)&=\frac{F_\Omega(t)}D
 =\frac1D\E_C\left[\sum_{j,k=1}^D e^{it(\theta_j-\theta_k)}\right],
 \qquad t\in\mathbb Z. \label{eq:k}
\end{align}
Aquí $e^{i\theta_j}$ son los eigenvalores de $U_C$. El promedio es sobre
monedas, no sobre estados iniciales. La invariancia Haar bajo
$C\mapsto e^{i\varphi}C$ implica
\begin{equation}
 \E_C\Tr U_C^t=0\quad(t\neq0).
\end{equation}
Por tanto, para tiempos enteros no nulos el SFF de este ensamble coincide
con el conectado. En $t=0$, $F_\Omega(0)=D^2$, $K_\Omega(0)=D$ y el
conectado es cero. Una densidad media uniforme del ensamble no obliga
a que la densidad de cada realización sea uniforme.

\section{Derivación de la referencia universal CUE}
Una moneda CUE(4) tiene sólo 16 parámetros reales; el operador de
dimensión $D$ no resulta Haar en $U(D)$ por esa sola elección.
La hipótesis adicional para usar CUE como clase universal es que
el operador completo tenga correlaciones caóticas de clase unitaria,
sin una simetría antiunitaria física que lo sitúe en COE, y que se
analice un espectro sin mezclar sectores conservados. Una moneda Haar
genérica no impone las relaciones de reciprocidad de una moneda COE.

Para una matriz \emph{completa} $V\sim\CUE(D)$, el núcleo de correlación
y la densidad de dos puntos distintos son
\begin{align}
 A_D(\theta-\theta')&=\frac1{2\pi}
        \sum_{m=0}^{D-1}e^{im(\theta-\theta')},\\
 \rho_2(\theta,\theta')&=\left(\frac D{2\pi}\right)^2
                  -|A_D(\theta-\theta')|^2.
\end{align}
Separando los $D$ términos diagonales de \eqref{eq:k}, para $t\neq0$,
\begin{equation}
 F_{\CUE(D)}(t)=D+int_0^{2\pi}\!\int_0^{2\pi}
  e^{it(\theta-\theta')}\rho_2(\theta,\theta')\,d\theta\,d\theta'.
\end{equation}
El término constante de $\rho_2$ integra a cero. La expansión del
cuadrado del núcleo deja una suma sobre $m,n\in\{0,\ldots,D-1\}$.
La integral sólo selecciona $t+m-n=0$: hay $D-|t|$ parejas para
$0<|t|<D$ y ninguna para $|t|\geq D$. Así,
\begin{equation}
 \boxed{F_{\CUE(D)}(t)=\min(|t|,D),\qquad
        K_{\CUE(D)}(t)=\min\left(\frac{|t|}{D},1\right)}
 \quad(t\neq0).
 \label{eq:cue}
\end{equation}
Este resultado finito es exacto para Haar en $U(D)$ \cite{bertini}.
Como $\Delta\theta=2\pi/D$, el tiempo de Heisenberg en pasos es
$t_H=2\pi/\Delta\theta=D=4N$. Con $\tau=t/t_H$, la predicción para
el caminante en el régimen universal es
\begin{equation}
 \boxed{K_\Omega(\tau)\simeq
 \begin{cases}\tau,&\tau_{\rm Th}\ll\tau\leq1,\\
               1,&\tau\geq1,
 \end{cases}\qquad \tau_{\rm Th}=t_{\rm Th}/D.}
\end{equation}
El tiempo de Thouless, las contribuciones de caminos cortos y la
dependencia de la densidad en cada moneda no quedan determinados por
la afirmación de caos geométrico. No se deduce una identidad a todos
los tiempos a partir de esa afirmación. La saturación del SFF de la
\emph{moneda} en $t=4$ tampoco determina la saturación del caminante.

\section{Resultados Haar exactos para el shift del repositorio}
\subsection{Un paso: contribución de frontera}
Sea $B_b$ el número de sitios cuyo vecino en la dirección $b$ no existe.
Los desplazamientos interiores no contribuyen a la traza. Sólo
contribuyen las reflexiones, por lo que
\begin{equation}
 \Tr U_C=\sum_{b=1}^4 B_b C_{b,\bar b}.
\end{equation}
Usando $\E[C_{ij}\overline{C_{kl}}]=\delta_{ik}\delta_{jl}/4$,
\begin{equation}
 \boxed{F_\Omega(1)=\frac14\sum_b B_b^2,\qquad
 K_\Omega(1)=\frac1{4D}\sum_b B_b^2
            =\frac1{16N}\sum_b B_b^2.}
 \label{eq:k1}
\end{equation}
Es un resultado exacto del modelo, válido sin asumir caos. Las
amplitudes se suman coherentemente sobre sitios porque todos comparten
la misma moneda. Para $L=49,R=37$, los parámetros del bloque espectral
de Sinaí en \texttt{FunctionsTester.nb}, se obtiene
\begin{equation}
 N=708,\quad D=2832,\quad (B_\uparrow,B_\downarrow,B_\rightarrow,B_\leftarrow)
 =(23,23,50,50),\quad K_\Omega(1)=0.534781\ldots .
\end{equation}
En cambio, la referencia Haar completa sería $1/2832=0.0003531\ldots$.
Esto prueba que \eqref{eq:cue} no describe exactamente el promedio
sobre monedas a todos los tiempos.

\subsection{Dos pasos: suma de caminos y cuarto momento Haar}
Definamos $B_{bd}$ como el número de sitios donde faltan ambos vecinos
$b$ y $d$. Los únicos caminos cerrados de dos pasos son un avance y
su regreso, o dos reflexiones en la misma posición:
\begin{equation}
 \Tr U_C^2=\sum_b (N-B_b)C_{b,\bar b}C_{\bar b,b}
       +\sum_{b,d}B_{bd}C_{b,\bar d}C_{d,\bar b}.
 \label{eq:trace2}
\end{equation}
Para dar una fórmula compacta y directamente implementable, escribamos
$\Tr U_C^2=\sum_{ijkl}A_{ijkl}C_{ij}C_{kl}$, donde $A$ cuenta los
caminos ordenados de \eqref{eq:trace2}. El cuarto momento Haar de $U(4)$ es
\begin{align}
 &\E[C_{ij}C_{kl}\overline{C_{ab}C_{cd}}]\nonumber\\
 &=\frac{\delta_{ia}\delta_{kc}\delta_{jb}\delta_{ld}
             +\delta_{ic}\delta_{ka}\delta_{jd}\delta_{lb}}{15}
 -\frac{\delta_{ia}\delta_{kc}\delta_{jd}\delta_{lb}
             +\delta_{ic}\delta_{ka}\delta_{jb}\delta_{ld}}{60}.
\end{align}
En consecuencia, con producto interno de Frobenius,
\begin{equation}
 \boxed{\begin{aligned}F_\Omega(2)&=
 \frac{\|A\|_F^2+\langle A,A^{\rm par}\rangle}{15}
 -\frac{\langle A,A^{\rm fila}\rangle+
         \langle A,A^{\rm col}\rangle}{60},\\
 K_\Omega(2)&=F_\Omega(2)/D.\end{aligned}}
\end{equation}
donde $A^{\rm par}_{ijkl}=A_{klij}$,
$A^{\rm fila}_{ijkl}=A_{kjil}$ y $A^{\rm col}_{ijkl}=A_{ilkj}$.
El script construye $A$ a partir de la permutación $S$ y evalúa esta
fórmula sin diagonalizar el operador.

En una familia bidimensional compacta con $B_b=O(\sqrt N)$,
el término interior dominante es
$2N(C_{\uparrow\downarrow}C_{\downarrow\uparrow}
    +C_{\rightarrow\leftarrow}C_{\leftarrow\rightarrow})$.
Como los dos monomios no tienen correlación Haar entre sí y cada uno
tiene segundo momento $1/15$,
\begin{equation}
 F_\Omega(2)=\frac{8N^2}{15}+O(N^{3/2}),\qquad
 K_\Omega(2)=\frac{2N}{15}+O(\sqrt N).
\end{equation}
La contribución temprana puede crecer con el tamaño aun cuando las
correlaciones locales posteriores sean caóticas.

\subsection{Estructura exacta a tiempos generales}
La traza es la suma de amplitudes de caminos cerrados con estados de
moneda: $\Tr U_C^t=\sum_{p\in\mathcal P_t}\prod_{r=1}^t
C_{b_r(p),a_r(p)}$. Sustituir en \eqref{eq:exactintegral} expresa el
SFF como suma sobre pares de caminos, seguida de una integral Haar
de $t$ entradas de $C$ y $t$ conjugadas. Esas integrales se evalúan
mediante contracciones de Weingarten. A partir de $t>4$, la matriz
Gram de permutaciones es singular y se usa su pseudoinversa o la
expansión restringida a particiones con a lo sumo cuatro filas.
No se puede factorizar el promedio por sitios ni por pasos: se repite
la misma variable aleatoria $C$. La geometría entra en las
multiplicidades de caminos; la hipótesis de caos permite una
aproximación universal, pero no fija esas multiplicidades exactas.

\section{Comprobación numérica y alcance}
Se generaron monedas Haar mediante QR de matrices Ginibre complejas
y corrección de las fases de la diagonal de $R$ \cite{mezzadri}.
Se diagonalizó $U_C$ y se calculó el SFF de las fases originales para
$t=1,\ldots,2D$. Estos tamaños son una comprobación a escala reducida,
no una diagonalización del estadio $L=120,R=90$ del PDF ni del caso
$L=49,R=37$ de los paneles espectrales.

\begin{center}
\begin{tabular}{lrr}
\toprule
 & $L=16,R=12$ & $L=24,R=18$\\
\midrule
Posiciones $N$ & 88 & 184\\
Dimensión $D$ & 352 & 736\\
Monedas independientes & 128 & 64\\
$K(1)$ exacto & 0.501420 & 0.522418\\
$K(1)$ numérico & $0.4774\pm0.0436$ & $0.4061\pm0.0593$\\
$K(2)$ exacto & 8.90398 & 20.09429\\
$K(2)$ numérico & $9.3879\pm1.3587$ & $16.5199\pm2.5082$\\
$\overline K$, $1.2\leq\tau\leq2$ & $0.9440\pm0.0066$ & $0.9420\pm0.0063$\\
\bottomrule
\end{tabular}
\end{center}
Los errores son errores estándar entre monedas. Para bandas temporales
se promedia primero dentro de cada moneda, conservando las
correlaciones entre tiempos al estimar el error. La figura
\texttt{sff\_cue\_2d.png} muestra bloques de aproximadamente $0.04D$
pasos y, aparte, los primeros 40 pasos.

Se observa una rampa y una aproximación a la meseta, con desviaciones
sistemáticas frente a la referencia CUE que superan el error estadístico.
No se ha demostrado convergencia exacta a CUE ni se ha determinado
$t_{\rm Th}$. La meseta sola no prueba caos: para cualquier espectro
simple, el promedio de Cesàro a tiempo infinito de $K$ es uno;
con degeneraciones de multiplicidad $m_\nu$ es
$\sum_\nu m_\nu^2/D$.

Como diagnóstico separado se desplegó cada espectro con una densidad
positiva suavizada por Fejér, usando 8 y 16 armónicos. Para $D=736$, la
media en $1.2\leq\tau\leq2$ pasa a $0.9764\pm0.0040$ y
$0.9838\pm0.0038$, respectivamente. El desplegado mejora parte del
acuerdo, pero deja desviaciones y dependencia del suavizado. Este
SFF de niveles desplegados no es $|\Tr U_C^t|^2/D$ del operador físico.
Los dos observables se guardan por separado.

La validación comprueba la unitariedad, que $S$ sea una permutación,
la igualdad entre trazas de potencias y eigenvalores, y los momentos
Haar de uno y dos pasos con 12000 monedas en una geometría pequeña.
Además, para una sola posición se recupera $F(1)=1,F(2)=2$ de CUE(4).
El máximo desvío del módulo de los eigenvalores en los dos ensambles
es $3.4\times10^{-14}$.

\section{Archivos y uso en Wolfram}
La carpeta de esta nota contiene \texttt{calcular\_sff.py}, los
eigenfases, monedas, geometría y series en \texttt{datos\_sff.npz},
tablas CSV, resúmenes JSON, la figura PNG/SVG y el diagnóstico de
desplegado. Las semillas son 20260930 y 20260931. El cálculo de
eigenfases y SFF requiere sólo NumPy; la figura usa Matplotlib.

En \texttt{QWMisc.wl}, la referencia de \texttt{qwSffPlot} usa
\texttt{SpectralFormFactorRMT[tau,1]} y la leyenda COE. Para CUE, el
índice debe ser 2. Los archivos originales se conservaron. Para un
SFF físico desde una lista completa de eigenfases, basta evaluar
\begin{verbatim}
dim = Length[eigenphases];
times = Range[2 dim];
k = (Abs[Total[Exp[I # eigenphases]]]^2/dim &) /@ times;
reference = Min[#, 1] & /@ N[times/dim];
\end{verbatim}
Si se usa \texttt{CompileSFF} sobre niveles desplegados, se está
calculando el segundo observable descrito arriba. En ambos casos la
dimensión para escalar el tiempo es el número de eigenfases, $D=4N$.

\begin{thebibliography}{9}
\bibitem{model} \emph{Discrete-time quantum walks in billiards},
\texttt{Chaos\_Results/main.pdf}, secciones I y III.D;
\texttt{QWMisc.wl}, \texttt{QuantumWalks/Billiards/Common.wl},
\texttt{Sinai.wl} y \texttt{FunctionsTester.nb} del repositorio local.
\bibitem{bertini} B. Bertini, P. Kos y T. Prosen,
\emph{Random Matrix Spectral Form Factor of Dual-Unitary Quantum Circuits},
secciones 2.3--2.4, ecuación (18),
\href{https://arxiv.org/abs/2012.12254}{arXiv:2012.12254}.
Se usa aquí su referencia CUE; no se traslada su teorema de circuitos
dual-unitarios al caminante.
\bibitem{mezzadri} F. Mezzadri,
\emph{How to generate random matrices from the classical compact groups},
\href{https://arxiv.org/abs/math-ph/0609050}{arXiv:math-ph/0609050}.
\end{thebibliography}
\end{document}
